\documentclass[11pt]{article}
\usepackage[margin=1.1in]{geometry}
\usepackage{amsmath,amssymb,amsthm}
\usepackage{booktabs}
\usepackage[hidelinks]{hyperref}
\usepackage{microtype}

\newcommand{\om}{\omega}
\newcommand{\vinf}{\|\om\|_\infty}
\newcommand{\Lam}{\Lambda}
\newcommand{\sL}{\sigma_{\!\Lambda}}
\newcommand{\Anaught}{A_0}

\title{The viscous inversion of the corner-flow blowup geometry fires at
the collapse transient's turnaround, at every viscosity tested}
\author{Justin Hill\\ \small Independent Research \\ \small justin@epagoge.io}
\date{Draft 0.1 --- 2026-08-02\thanks{Sequel to [P1]. Every number
regenerates from bytes on disk (Section~\ref{sec:repro}). This paper was
drafted after, not before, a two-lens verification pass over its own
claim records: roughly 230 numeric claims recomputed independently,
twelve scoping corrections applied, and four of our own estimators and
readings killed en route. The corpses are in Section~\ref{sec:corpses}.}}

\begin{document}
\maketitle

\begin{abstract}
Paper [P1] measured $\sL$, the scaling exponent of the scale-invariant
direction-regularity observable $\Lam=\sup_P|\nabla\xi|\,|\om|^{-1/2}$,
on the Boussinesq corner-flow blowup scenario: $+1.00\pm0.03$ inviscid
(calibrated where blowup is proven), inverting to a deep-collapse value
near $-1.2$ under viscosity at $\nu\in\{10^{-4},10^{-3}\}$. This paper
asks the questions that result forces: how the inversion depends on
viscosity, where in the collapse it happens, and what sets that
location. Within the tested scenario family (Section~\ref{sec:scope}):
(1)~\emph{No resolved viscosity dependence.} At the five tested
viscosities $\nu\in\{10^{-5},3\times10^{-5},10^{-4},3\times10^{-4},
10^{-3}\}$ the deep-collapse exponent shows no trend (chord estimator
$-1.12$ to $-1.29$, every cross-grid spread under the pre-registered
$0.15$ bar; two-slope-form asymptote $-1.07\pm0.03$); the two
previously certified viscosities replicated 4/4 against independent
runs. Against the window-matched inviscid anchor ($+0.589/+0.585$) the
data show no approach toward the inviscid value down to $\nu=10^{-5}$:
any crossover in $\nu$ lies below the tested range. The data are
consistent with a $\nu\to0^{+}$ discontinuity of the exponent; they
cannot distinguish it from a crossover below $10^{-5}$.
(2)~\emph{A fixed trigger.} The switch sits at
$A^{*}=52.7$--$54.2\times\Anaught$ in twelve measurements at
$\nu\in\{10^{-4},10^{-3}\}\times\Anaught\in\{50,100,200\}$, both grids,
and is bracketed inside the same band at the three other viscosities
($\Anaught=100$ only). Since the initial corner-gradient amplitude is
$37.7\times\Anaught$ (a geometric constant of this IC family), the
trigger is a growth factor $1.43$ (spread $1.40$--$1.44$) above the
initial value: about a third of an e-fold. The switch is nearly a step,
completing within ${\sim}0.2$--$0.3$ e-folds (measured at $\nu=10^{-4}$
only). (3)~\emph{The two instruments meet.} An empirical clock fit
($\gamma_A=0.77$--$1.11$) places the switch at $s^{*}=0.52\pm0.05$ in
rescaled time (per-run degeneracy widens this to $[0.29,0.70]$),
consistent with the corner-frame march's fast-transient turnaround at
$s=0.5$ (resolved to the $ds=0.25$ step); the pre-registered lock-on
hypothesis ($s\sim3.7$--$13.5$) is refuted by more than a factor of
five. A large-deviation march closes the linearity gap as far as the
march reaches: the $s=0.5$ turnaround is amplitude-invariant across
4.5 decades of perturbation amplitude, and the march's
basin-or-instrument edge lies between amplitudes $3\times10^{-2}$ and
$10^{-1}$ along the one direction tested. The licensed mechanism
statement, exactly bounded: viscosity is necessary for the inversion
(inviscid runs never invert), and at the two viscosities where the
clock was fit the inversion occurs at a rescaled time consistent with
the transient's turnaround, well before self-similar settling; the
switch location tracks a $\nu$-independent, $\Anaught$-linear stage of
the collapse whose dynamical identity remains open.
\end{abstract}

\section{Scope, stated first}\label{sec:scope}

Everything in this paper is measured on one scenario family: the
axisymmetric Boussinesq corner flow of the Luo--Hou type, quartic
initial condition, corner scale $r_0=0.4$, \texttt{zpow}~1, free-slip
viscosity in 5D-Laplacian form (the no-slip results of Barker--Prange
\cite{BP20} concern a different boundary condition and are cited as
prior art for the criteria, not as the same setting). Amplitudes are in
code units of the vector-$|\om|$ maximum; grids are $128\times384$ and
$256\times768$ (one refinement doubling; the $0.15$ cross-grid bar
tests consistency under that doubling, not asymptotic grid
convergence); the march runs seed-0 perturbations at $ds=0.25$ with the
M2 stop rule, and every accepted march step converged the true
nonlinear residual below $10^{-10}$.

\emph{Novelty scoping.} The prior-art search of [P1] (a one-session web
sweep; ``not found by this search,'' never ``proven absent'') covers
the viscosity-resolved exponent curve: no measurement of a
direction-regularity exponent as a function of viscosity was found, and
Section~\ref{sec:curve} is that curve. The trigger invariant, switch
width, clock comparison, and basin bound are reported as measurements
with no novelty claim attached. Adjacent art sighted in a supplementary
search: analytical viscous geometric depletion \cite{Ju06,Gru26}, and
nearly self-similar viscous blowup in \emph{generalized} axisymmetric
Navier--Stokes with dimension as a parameter \cite{Hou24}; our
measurements concern the standard equations at fixed dimension, on this
scenario family, at the viscosities stated.

\section{The curve: no resolved \texorpdfstring{$\nu$}{nu}-dependence}
\label{sec:curve}

Ten new DNS runs (five viscosities, both grids, doubled snapshot
cadence, fresh tags) went through the certified estimator of [P1]
verbatim: signed circulation-drift gate, peak-box $\Lam$ recipe,
cross-grid common-amplitude overlap, top half, symmetric midpoint cut,
OLS slope with $10^4$-resample bootstrap intervals.

\begin{center}
\begin{tabular}{lccc}
\toprule
$\nu$ & $128\times384$ & $256\times768$ & spread \\
\midrule
$0$ (anchor) & $+0.589$ & $+0.585$ & $0.005$ \\
$10^{-5}$ & $-1.174$ & $-1.190$ & $0.016$ \\
$3\times10^{-5}$ & $-1.271$ & $-1.156$ & $0.115$ \\
$10^{-4}$ & $-1.119$ & $-1.187$ & $0.068$ \\
$3\times10^{-4}$ & $-1.286$ & $-1.204$ & $0.082$ \\
$10^{-3}$ & $-1.207$ & $-1.231$ & $0.024$ \\
\bottomrule
\end{tabular}
\end{center}

Every spread passes the pre-registered $0.15$ bar. The two viscosities
certified in [P1] replicated 4/4 against these independent runs (worst
$|\Delta|=0.055$, within tolerance); the three new viscosities are
single run-pairs with cross-grid agreement as their check. The deep
windows carry $n=20$--$27$ snapshots (the doubled cadence fixed [P1]'s
thinnest tables). A two-slope-form fit (Section~\ref{sec:trigger}) puts
the asymptotic deep slope at $-1.07\pm0.03$ across all six fitted runs,
systematically shallower than the chord values; both are quoted and
neither replaces the other. The chord-vs-form gap (${\sim}0.1$) is the
current estimator systematic on the exponent's absolute value; the
$\nu$-flatness conclusion is insensitive to it.

What the table does not show: any approach toward the inviscid anchor.
If $\sL^{\mathrm{deep}}(\nu)$ returns to $+0.59$ as $\nu\to0$, it does
so below $\nu=10^{-5}$. Within the tested range the inversion is
viscosity-blind. The Corollary-0 reading of [P1] (the viscous cap on
direction-gradient growth exists for every $\nu>0$ and is absent at
$\nu=0$ \cite{Con94}) is consistent with a genuine $0^{+}$
discontinuity; that reading is interpretation, and the measurement
alone cannot exclude a crossover below the tested range.

\section{The trigger: a fixed growth factor, and a near-step switch}
\label{sec:trigger}

At finite $\nu$, $\sigma(A)$ is not one number: every viscous run
crosses from inviscid-like slope ($+0.39$ to $+0.81$) at low amplitude
to the inverted value at high amplitude. [P1] adjudicated
window-averaged slopes as regime mixtures; here the crossover itself
becomes the observable.

The estimator lineage is part of the record. A pre-registered
chord-intersection estimator for the crossover amplitude \emph{failed}
its split-invariance test (15--42\% drift tracking the window split,
monotone in nine of ten runs) and its outputs are void: $\sigma(A)$
curves smoothly rather than breaking, so a two-chord intersection
follows the split point. Its successor, the first zero crossing of a
$k$-row sliding-window local slope, passes the same class of test that
killed its predecessor: $k$-drift 0.9--2.1\% on the discovery runs and
0.7--1.0\% on the cross-term runs, cross-grid agreement under 1\%.

By that estimator, the crossover amplitude is $\Anaught$-linear with
measured power $+1.00$ (ratios $0.50/0.49$ at $\Anaught=50$ and
$2.01/2.00$ at $\Anaught=200$ against the amplitude-symmetry
prediction), and the twelve measurements at
$\nu\in\{10^{-4},10^{-3}\}\times\Anaught\in\{50,100,200\}$ give
$A^{*}/\Anaught=52.7$--$54.2$. At $\nu\in\{3\times10^{-4},
3\times10^{-5},10^{-5}\}$ ($\Anaught=100$ only) the crossover is
bracketed inside the same fixed band: the lower half of every overlap
window reads inviscid-positive and the upper half reads inverted, so
the crossover cannot have moved with $\nu$ within the tested window.
The tested region of the $(\nu,\Anaught)$ plane is a cross, not a tile;
the plane statement is exact for that cross.

The initial corner-gradient amplitude is $37.7\times\Anaught$ in every
run, a geometric constant of the quartic IC family (early inviscid
snapshots grow $37.71\to37.80\to37.93$: the value at the first trusted
snapshot is essentially the initial value). The trigger is therefore a
growth factor of $1.43$ above the initial amplitude (twelve-run spread
$1.40$--$1.44$; the tanh-form estimator places the crossing
systematically ${\sim}8\%$ lower, so the absolute factor is
estimator-conditional at that level while $\Anaught$-linearity and
$\nu$-independence are estimator-robust). Three dependencies are
declared: the factor is a property of this IC family, of the
vector-$|\om|$ observable, and of the zero-crossing estimator.

The switch is nearly a step. A two-slope blend fit is
residual-adequate on all six runs at $\nu=10^{-4}$ but
parameter-degenerate (the optimizer drives it into an
exponential-transient limit; its center, width, and low-$A$ slope are
not individually physical and are not quoted). The robust extract is
the transient decay scale: the local slope completes its transition
from $0$ to the inverted value within ${\sim}0.2$--$0.3$ e-folds of
amplitude. The width is measured at $\nu=10^{-4}$ only.

One confound was tested and killed before these claims were written:
across the switch, the $|\om|$ argmax stays at the corner ($z^{*}=0$,
$r^{*}=1.000$) at every gated snapshot, and the observable tracks the
swirl-gradient structure throughout ($\om_\theta$ vanishes at the
corner by symmetry). The diagnostic was run at $\nu=10^{-4}$ on the
128-grid and on the matched inviscid run; we treat it as excluding a
structure handoff for the runs where the switch is certified. The
inversion is a change in the direction-field geometry at one fixed
structure, not the instrument changing targets.

\section{The clock: two instruments meet at the transient's turnaround}
\label{sec:clock}

[P1]'s rescaled-frame march certified the corner profile as an
attractor: perturbations peak at $s=0.5$ (raw factor $2.389$ at linear
amplitude) and contract at $0.270$ per $s$-unit. The DNS switch can be
placed on the same $s$-axis by fitting the collapse clock
$\ln A=-\gamma_A\ln(1-t/T^{*})+\ln A_i$ per run, $T^{*}$ scanned, with
$s=-\ln(1-t/T^{*})$.

\begin{center}
\begin{tabular}{lccc}
\toprule
run & $\gamma_A$ & $s^{*}$ & $s^{*}$ over $2\times$-residual $T^{*}$ set \\
\midrule
OR $128\times384$ ($\nu=0$) & $0.77$ & $0.61$ & $[0.42, 0.70]$ \\
OR $256\times768$ ($\nu=0$) & $0.94$ & $0.54$ & $[0.41, 0.58]$ \\
C $128\times384$ ($10^{-4}$) & $0.86$ & $0.55$ & $[0.38, 0.61]$ \\
C $256\times768$ ($10^{-4}$) & $1.11$ & $0.48$ & $[0.35, 0.51]$ \\
C $128\times384$ ($10^{-3}$) & $1.04$ & $0.47$ & $[0.30, 0.53]$ \\
C $256\times768$ ($10^{-3}$) & $1.03$ & $0.49$ & $[0.31, 0.55]$ \\
\bottomrule
\end{tabular}
\end{center}

$s^{*}=0.52\pm0.05$ (run-to-run scatter; the $T^{*}$--$\gamma_A$
degeneracy widens the honest per-run interval to $[0.29,0.70]$),
consistent with the march turnaround at $s=0.5$, itself resolved only
to the $ds=0.25$ step. The value is conditional on the clock form and
on the $54\times\Anaught$ switch input: the tanh-implied
$49$--$50\times\Anaught$ shifts $s^{*}$ by about $-0.09$, inside the
degeneracy band. $\gamma_A$ spans $0.77$--$1.11$, consistent with the
corner frame's one-e-folding-per-$s$-unit and measured rather than
assumed. The pre-registered hypothesis space put lock-on at
$s\sim3.7$--$13.5$; even the widest degeneracy bound clears that window
by more than a factor of five. Within the pre-registered space and
clock family, the lock-on hypothesis is refuted: the switch does not
wait for the collapse to settle onto the profile.

The comparison identifies matching timescales of one operator, not
coincident trajectories: the march transient is measured for
perturbations about the settled profile, while the DNS approach is a
large deviation. The next section closes that gap as far as the march
can reach.

\section{The large-deviation march}\label{sec:ld}

The certified march ran the same seed-0 perturbation at amplitudes
$3\times10^{-3}$, $10^{-2}$, $3\times10^{-2}$, and $10^{-1}$ (the
largest previously validated amplitude was $10^{-3}$). At the first
three, all 240 steps are valid and the raw-norm peak sits at
$s_{\mathrm{peak}}=0.5$ exactly, at step resolution. Combined with
[P1]'s M2 ladder, the transient turnaround is amplitude-invariant
across 4.5 decades of relative deviation ($10^{-6}$ to
$3\times10^{-2}$). The $s=0.5$ timescale that the DNS switch lands on
is a property of the operator's nonlinear neighborhood, not merely of
its linearization.

Around the pinned turnaround, genuine nonlinear structure appears for
the first time in this campaign: the transient's end stretches
($s_{\mathrm{transient}}=13.5\to14.0\to14.75$, $+9\%$) and the peak
factor grows ($2.4068\to2.4504\to2.5884$ within this ladder; M2's
linear value is $2.389$) while the peak position does not move. The
Lyapunov $P$-norm contracts monotonically from $V_0=77$ through 13.7
decades with zero violations above the pre-registered $V$-floor
artifact threshold: the certificate's contraction holds at 3\%
deviation.

At amplitude $10^{-1}$ the march stalls at step one (raw growth
$3.3\times$ in one step; the $P$-norm nevertheless decreased on that
step, a mixed signal). The pre-registered protocol does not separate
dynamical escape from solver stall, so the honest statement is a
basin-or-instrument edge: along the seed-0 direction,
$3\times10^{-2}<r<10^{-1}$ in the march's norm. This also bounds the
remaining daylight in Section~\ref{sec:clock}'s comparison: the DNS
initial deviation exceeds this basin, and its approach enters the
basin during the transient. Full closure would need DNS-state
initialization of the march, which is a different instrument build.

\section{Mechanism, stated within bounds}\label{sec:mechanism}

For every viscosity tested, the inversion is present only with
$\nu>0$: inviscid runs on identical amplitude windows never invert. At
the two viscosities where the clock was fit, the switch occurs at a
rescaled time consistent with the corner-frame transient's turnaround,
well before self-similar settling, and the turnaround timescale is
amplitude-robust across every amplitude the march can hold. Viscosity
is necessary for the inversion; the location of the switch tracks a
$\nu$-independent, $\Anaught$-linear stage of the collapse. What sets
that stage dynamically is open, and Section~\ref{sec:curve}'s gap below
$\nu=10^{-5}$ is open with it. The strongest reading consistent with
all measurements: the geometric depletion that [P1] measured deep in
the collapse is already fully switched on by the end of the approach
transient, at every viscosity and amplitude tested, and the
deep-collapse exponent it switches to does not depend on how much
viscosity did the switching.

\section{Relation to prior art}\label{sec:prior}

The criteria territory is [P1]'s: Constantin--Fefferman \cite{CF93},
Beir\~ao da Veiga--Berselli \cite{BdVB02,BdVB09} at the definition of
$\Lam$, Giga--Miura \cite{GM11} and Barker--Prange \cite{BP20} for
qualitative type-I exclusion under direction coherence, and
Constantin's identity \cite{Con94} behind the Corollary-0 diagnostic.
Analytical viscous geometric depletion is developed in \cite{Ju06} and,
through logarithmically weighted direction classes, in \cite{Gru26};
neither defines a measured exponent or an onset. Nearly self-similar
viscous blowup is reported in \emph{generalized} axisymmetric
Navier--Stokes with dimension as a parameter \cite{Hou24}; the present
measurements concern the standard equations on this scenario family
and are not in tension with results for modified systems. The measured
objects of this paper (an exponent-vs-viscosity curve, a trigger
invariant, a cross-instrument clock comparison) were not found in the
[P1] search; only the curve carries a novelty claim, under that
search's stated horizon.

\section{What died on the way}\label{sec:corpses}

Four of our own results were killed by our own instruments during this
campaign, and the corpses are retained:

\begin{itemize}
\item The chord-intersection crossover estimator: failed
  split-invariance (15--42\% drift); all its outputs, including a
  fitted crossover-vs-viscosity slope, are void.
\item The tanh-form parameters: residual-adequate fit, degenerate
  parameters; only the decay scale and asymptote survive.
\item The ``four e-folds of dynamical age'' reading of the trigger:
  killed by the normalization audit (the initial corner-gradient
  amplitude is $37.7\times\Anaught$, not $\Anaught$); the honest
  trigger is $1.43\times$ above the initial value, essentially at
  collapse onset.
\item The lock-on bridge hypothesis: the clocks were compared and it
  lost, by a factor of five beyond the widest degeneracy bound.
\end{itemize}

[P1]'s standing retractions (window-averaged viscous slopes; the
amplitude-composite reading of the full-window inviscid calibration)
are inherited: no full-window $\sigma$ is quoted anywhere in this
paper, and every viscous contrast is drawn against the window-matched
inviscid $+0.589$, never against $+1.00$.

\section{Reproducibility}\label{sec:repro}

The DNS campaign is 18 runs (about 90 minutes of laptop wall time)
plus the march ladders (${\sim}90$ seconds each), all regenerable: run
scripts, pre-registration (every estimator, bar, and hypothesis space
in this paper was written down before the data it judges existed; the
spec file's section dates interleave with the run logs), extraction
code, and per-claim records with their adjudication addenda are in the
repository alongside [P1]'s verification suite (15/15 standing,
watcher live). The two-lens audit that preceded this draft (a numbers
tracer recomputing ${\sim}230$ claims from the raw arrays and
snapshots; an adversarial referee producing the twelve scoping
corrections now baked into the text) is itself part of the record.


\section{Tool and computational resource disclosure}\label{sec:disclosure}

As in [P1], and following the Leiden Declaration \cite{Leiden26}:

\textbf{Tools.} Dedalus for the simulations, NumPy/SciPy for the
estimators, no proof assistant. An AI assistant (Anthropic Claude, used
interactively) wrote the run drivers, the tracking and estimator code,
and this paper's prose, and executed the supplementary prior-art search
of Section~\ref{sec:prior}.

\textbf{Division of labour.} The code and prose are machine-assisted;
the protocol is the author's, and in this paper the protocol is most of
the method: every estimator, bar, and hypothesis space in
Sections~\ref{sec:curve}--\ref{sec:ld} was written down before the data
it judges existed (the specification file's section timestamps
interleave with the run logs), and two pre-registered routes were
stopped by their own tests rather than rescued. The four corpses in
Section~\ref{sec:corpses} are machine-produced results that the author's
protocol killed.

\textbf{What the author can and cannot defend at expert level.} This
paper is almost entirely measurement, and its exposure is
correspondingly narrower than [P1]'s. The author can defend the
estimators, their invariance tests, the pre-registration record, and
every retraction. The one inherited analytical dependency is [P1]'s
exclusion chain, which is cited here for context and carries the same
caveat stated in [P1]'s disclosure: it is machine-derived, conditional
on a structural hypothesis with a measured constant, and unverified.
No measurement in this paper depends on it.

\textbf{Responsibility.} Rests with the human author, including for the
machine-assisted portions.

\begin{thebibliography}{99}

\bibitem{Leiden26} Leiden Declaration on Artificial Intelligence
and Mathematics. 2 June 2026. doi:10.5281/zenodo.20302944.
\url{https://leidendeclaration.ai}


\bibitem{P1} J.~Hill.
Viscosity inverts the direction-regularity scaling of the corner-flow
blowup mechanism: a measured exponent crosses the type-I exclusion
threshold. 2026.
\url{https://epagoge.github.io/parzival/sigma/}
(paper, LaTeX, data, and claims ledger; repository EPAGOGE/parzival).

\bibitem{CF93} P.~Constantin, C.~Fefferman.
Direction of vorticity and the problem of global regularity for the
Navier--Stokes equations.
\emph{Indiana Univ. Math. J.} \textbf{42}, 775--789 (1993).

\bibitem{Con94} P.~Constantin.
Geometric statistics in turbulence.
\emph{SIAM Review} \textbf{36}, 73--98 (1994).

\bibitem{BdVB02} H.~Beir\~ao da Veiga, L.C.~Berselli.
On the regularizing effect of the vorticity direction in
incompressible viscous flows.
\emph{Differential Integral Equations} \textbf{15} (2002).

\bibitem{BdVB09} H.~Beir\~ao da Veiga, L.C.~Berselli.
Navier--Stokes equations: Green's matrices, vorticity direction, and
regularity up to the boundary (2009).

\bibitem{GM11} Y.~Giga, H.~Miura.
On vorticity directions near singularities for the Navier--Stokes
flows with infinite energy.
\emph{Comm. Math. Phys.} (2011).

\bibitem{BP20} T.~Barker, C.~Prange.
Scale-invariant estimates and vorticity alignment for Navier--Stokes
in the half-space with no-slip boundary conditions.
\emph{Arch. Ration. Mech. Anal.} \textbf{235}, 881--926 (2020).
arXiv:1906.08225.

\bibitem{Ju06} N.~Ju.
Geometric depletion of vortex stretch in 3D viscous incompressible
flow.
\emph{J. Math. Anal. Appl.} \textbf{321}(1), 412--425 (2006).
doi:10.1016/j.jmaa.2005.08.048.

\bibitem{Gru26} Z.~Gruji\'c.
Logarithmic depletion of vortex stretching and singularity evasion in
the 3D Navier--Stokes equations.
arXiv:2607.08866 (2026).

\bibitem{Hou24} T.Y.~Hou.
Nearly self-similar blowup of generalized axisymmetric Navier--Stokes
equations.
arXiv:2405.10916 (2024, revised 2025).
[Verified 2026-08-02: generalized equations, dimension as parameter;
preprint.]

\end{thebibliography}

\end{document}
